\subsection{The exact meaning of vanishing odd torsion}
\label{subsec:odd-torsion-plus-adjoint}

The odd torsion $D$ is related to the plus class module
by the following exact adjoint comparison.

\begin{proposition}\label{prop:D-plus-finiteness}
For every odd prime $p$,
\[
                             D=0\quad\Longleftrightarrow\quad X^+\text{ is finite}.
\]
Under these equivalent assertions the proved
single-component theorem also gives $Z=0$.
\end{proposition}

\begin{proof}
For an odd exponent $a\ne1,-1$, set
$\chi=\omega^{-a}$ and $\theta=\omega^{1-a}$.
The proof of \cite[Proposition~2.2.6, p.~4]{WWEOrd}
gives the actual Poitou--Tate identification
\[
 \operatorname{Ext}^1_\Lambda
       ((M^{[a]})^\#(1),\Lambda)=X^{[1-a]}.
\]
The involution and twist act on the cyclotomic
variable by $\gamma\mapsto u\gamma^{-1}$.
At tame level one, the source assumes $p\geq5$,
$\theta$ primitive and $\theta\ne\omega^2$. These
conditions exclude precisely $a=1,-1$, which we treat below.

Put $R=\Lambda[1/p]$. Restricting extensions along
$D^{[a]}\hookrightarrow M^{[a]}$, after the same
twist and localization, gives an isomorphism
\[
 X^{[1-a]}[1/p]\xrightarrow{\ \sim\ }
 \operatorname{Ext}^1_R
       ((D^{[a]})^\#(1)[1/p],R).
\]
Indeed the torsionfree quotient is projective over
the one-dimensional PID $R$, so the long exact
Ext sequence proves this exact restriction map is
an isomorphism. A nonzero cyclic torsion summand
$R/(f^k)$ has Ext group $R/(f^k)$, by its length-one
resolution. Consequently $D^{[a]}=0$ exactly when
$X^{[1-a]}[1/p]=0$. The first module has no
$\mathbf Z_p$-torsion; Ferrero--Washington makes
the second finitely generated over $\mathbf Z_p$,
so its rational vanishing is precisely finiteness.

For the exceptional components, norm to $\mathbf Q$
and extension back give $A_1^{[0]}=0$, since their
composite is multiplication by the unit $p-1$.
Herbrand gives $A_1^{[1]}=0$ and, for $p>3$,
$A_1^{[-1]}=0$ from $B_2=1/6$; see also
\cite[Lemma~5.2.7, p.~13]{WWEOrd}.
Leopoldt reflection with $\chi=\omega^{-1}$ gives
$r_p(A_1^{[2]})\leq r_p(A_1^{[-1]})=0$.
These are the exact cases of
\cite[Theorems~4.4.2 and~4.5.7]{SharifiIwa}.
Proposition~\ref{prop:component-control} therefore
gives $X^{[0]}=X^{[2]}=X^{[1]}=X^{[-1]}=0$.

The class-field unit sequence and the fact that the
odd global units are exactly $\mathbf Z_p(1)$,
together with Coleman's exact sequence
\cite[Theorem~5.5.31, p.~148]{SharifiIwa}, give the
actual maps
\[
 M^{[1]}\simeq\ker(\Lambda\xrightarrow{\gamma\mapsto u}
                \mathbf Z_p(1))=(\gamma-u)\Lambda,
 \qquad M^{[-1]}\simeq\Lambda\quad(p>3).
\]
Hence both exceptional odd torsion components vanish.
At $p=3$ these odd characters coincide, as do the
even characters $0,2$; the first displayed map and
the trivial-character class calculation cover them.
The correspondence $a\mapsto1-a$ now covers every
odd and every even character, proving the equivalence.
Applying the established finite-$X^{[e]}$ component
theorem to every even $e$ gives $Z=0$.
\end{proof}

The exact adjoint map therefore identifies the assertion
$D=Z=0$ with the finite-$X^+$ hypothesis. The canonical
finite-stage quotients are defined at every odd prime,
including those for which this hypothesis remains undecided.
